\[ %%%%%%%%%%%%%%%%%%%%%%%%%%%%%%%%%%%%%%%%%%%%%%%%%%%%%%%%%%%%%%%%%%%%%%%%%%%%%%%% \subsection{Dynamical System Formulation} \label{sec:dynamical_system} %%%%%%%%%%%%%%%%%%%%%%%%%%%%%%%%%%%%%%%%%%%%%%%%%%%%%%%%%%%%%%%%%%%%%%%%%%%%%%%% Physical systems considered in this work are formulated as nonlinear dynamical systems whose evolution is governed by underlying physical laws. Depending on the application domain, these systems may be described by ordinary differential equations (ODEs), partial differential equations (PDEs), or discrete transition models. The general continuous-time formulation of a physical system is: \begin{equation} \frac{dX(t)}{dt} = F(X(t),t,\Theta), \end{equation} where: \begin{itemize} \item $X(t)$ represents the physical state of the system; \item $F(\cdot)$ represents the unknown or partially known physical evolution operator; \item $\Theta$ represents physical parameters and system properties. \end{itemize} The objective of PHYSGEN is to construct a learnable approximation: \begin{equation} F_\theta(X,t) \approx F(X,t,\Theta), \end{equation} capable of predicting future system states while preserving physical structure. %%%%%%%%%%%%%%%%%%%%%%%%%%%%%%%%%%%%%%%%%%%%%%%%%%%%%%%%%%%%%%%%%%%%%%%%%%%%%%%% \subsubsection{Discrete-Time Dynamical Representation} %%%%%%%%%%%%%%%%%%%%%%%%%%%%%%%%%%%%%%%%%%%%%%%%%%%%%%%%%%%%%%%%%%%%%%%%%%%%%%%% For computational simulation and machine learning applications, continuous dynamics are typically discretized in time: \begin{equation} X_{t+\Delta t} = X_t + \Delta t F(X_t,\Theta). \end{equation} This formulation defines the physical transition operator: \begin{equation} X_{t+1} = \mathcal{T}(X_t), \end{equation} where $\mathcal{T}$ represents the discrete evolution mapping. PHYSGEN learns an approximation of this transition: \begin{equation} \hat{X}_{t+1} = \mathcal{T}_{\theta}(X_t). \end{equation} However, unlike conventional neural transition models, PHYSGEN represents the system state through a dynamically evolving interaction graph. %%%%%%%%%%%%%%%%%%%%%%%%%%%%%%%%%%%%%%%%%%%%%%%%%%%%%%%%%%%%%%%%%%%%%%%%%%%%%%%% \subsubsection{Graph Dynamical System Formulation} %%%%%%%%%%%%%%%%%%%%%%%%%%%%%%%%%%%%%%%%%%%%%%%%%%%%%%%%%%%%%%%%%%%%%%%%%%%%%%%% The physical state is transformed into a graph representation: \begin{equation} X_t \rightarrow G_t=(V_t,E_t,X_t,P_t). \end{equation} The evolution of the physical system is therefore reformulated as a graph dynamical process: \begin{equation} G_{t+1} = \Phi(G_t,P_t). \end{equation} PHYSGEN learns the approximation: \begin{equation} \hat{G}_{t+1} = \Phi_{\theta}(G_t,P_t). \end{equation} The learned operator consists of several coupled transformations: \begin{equation} \Phi_{\theta} = \Phi_{dec} \circ \Phi_{dyn} \circ \Phi_{msg} \circ \Phi_{enc}, \end{equation} where: \begin{itemize} \item $\Phi_{enc}$ maps physical states into latent graph representations; \item $\Phi_{msg}$ performs physics-guided information propagation; \item $\Phi_{dyn}$ models latent physical evolution; \item $\Phi_{dec}$ reconstructs future physical states. \end{itemize} %%%%%%%%%%%%%%%%%%%%%%%%%%%%%%%%%%%%%%%%%%%%%%%%%%%%%%%%%%%%%%%%%%%%%%%%%%%%%%%% \subsubsection{Connection with Differential Equations} %%%%%%%%%%%%%%%%%%%%%%%%%%%%%%%%%%%%%%%%%%%%%%%%%%%%%%%%%%%%%%%%%%%%%%%%%%%%%%%% Many physical systems can be represented as spatially distributed fields: \begin{equation} u(x,t), \end{equation} whose evolution is governed by: \begin{equation} \frac{\partial u}{\partial t} = \mathcal{F} ( u, \nabla u, \nabla^2u, \Theta ). \end{equation} Examples include: \begin{itemize} \item Navier--Stokes equations for fluid dynamics; \item heat diffusion equations; \item reaction--diffusion systems; \item electromagnetic field equations. \end{itemize} After spatial discretization, these continuous fields become collections of interacting local states: \begin{equation} u_i^{t+1} = u_i^t + \Delta t \, F_i ( u_i^t, \{u_j^t:j\in\mathcal{N}(i)\} ). \end{equation} PHYSGEN generalizes this local evolution principle by replacing the manually defined operator $F_i$ with a physics-guided graph evolution operator: \begin{equation} h_i^{t+1} = \Phi_\theta ( h_i^t, \{h_j^t:j\in\mathcal{N}(i)\}, P_i ). \end{equation} %%%%%%%%%%%%%%%%%%%%%%%%%%%%%%%%%%%%%%%%%%%%%%%%%%%%%%%%%%%%%%%%%%%%%%%%%%%%%%%% \subsubsection{State Space and Physical Manifold} %%%%%%%%%%%%%%%%%%%%%%%%%%%%%%%%%%%%%%%%%%%%%%%%%%%%%%%%%%%%%%%%%%%%%%%%%%%%%%%% A physical dynamical system does not evolve arbitrarily in the entire state space. Instead, physically valid trajectories occupy a constrained subset: \begin{equation} X(t)\in\mathcal{M}_{phys}, \end{equation} where $\mathcal{M}_{phys}$ denotes the physical state manifold. The objective of PHYSGEN is therefore formulated as learning an evolution operator satisfying: \begin{equation} \Phi_\theta: \mathcal{M}_{phys} \rightarrow \mathcal{M}_{phys}. \end{equation} This condition distinguishes physically consistent prediction from purely statistical forecasting, where predicted trajectories may leave the admissible region of the system dynamics. %%%%%%%%%%%%%%%%%%%%%%%%%%%%%%%%%%%%%%%%%%%%%%%%%%%%%%%%%%%%%%%%%%%%%%%%%%%%%%%% \subsubsection{Nonlinear Long-Horizon Evolution} %%%%%%%%%%%%%%%%%%%%%%%%%%%%%%%%%%%%%%%%%%%%%%%%%%%%%%%%%%%%%%%%%%%%%%%%%%%%%%%% For long-term forecasting, the learned operator is recursively applied: \begin{equation} \hat{X}_{t+k} = \Phi_\theta^k(X_t). \end{equation} The main challenge is that small approximation errors may accumulate: \begin{equation} \epsilon_k = \| X_{t+k} - \Phi_\theta^k(X_t) \|. \end{equation} PHYSGEN addresses this problem through the combination of: \begin{itemize} \item physics-guided message propagation; \item latent physical state evolution; \item adaptive stability control; \item multi-level physical constraints. \end{itemize} %%%%%%%%%%%%%%%%%%%%%%%%%%%%%%%%%%%%%%%%%%%%%%%%%%%%%%%%%%%%%%%%%%%%%%%%%%%%%%%% \subsubsection{Summary} %%%%%%%%%%%%%%%%%%%%%%%%%%%%%%%%%%%%%%%%%%%%%%%%%%%%%%%%%%%%%%%%%%%%%%%%%%%%%%%% The dynamical system formulation establishes the mathematical foundation of PHYSGEN as a learnable approximation of nonlinear physical evolution operators. By transforming continuous physical dynamics into adaptive graph-based representations, PHYSGEN creates a unified framework connecting differential equations, graph neural networks, and physics-guided machine learning. \] 
